\question{}

این ادعا به طور کلی نادرست است و می‌توانیم برای آن مثال نقض بیاوریم.

فرض کنید در یک \lr{MDP}، مجموعه‌ی حالت‌ها $\{s_0, s_1, s_2, s_3\}$ باشد که $s_3$ حالت پایانی است، و مجموعه‌ی اکشن‌ها هم $\{a_0, a_1\}$ باشد. همچنین برای سادگی فرض می‌کنیم که ضریب تخفیف نداشته باشیم (یا به طور معادل $\gamma = 1$). همچنین توابع انتقال و پاداش‌ها به این صورت هستند:
\begin{itemize}
    \item \textbf{از حالت $s_0$:}
          \begin{itemize}
              \item با انتخاب اکشن $a_0$ قطعاً به حالت $s_1$ می‌رویم (پاداش = $0$).
              \item با انتخاب اکشن $a_1$ قطعاً به حالت $s_2$ می‌رویم (پاداش = $0$).
          \end{itemize}

    \item \textbf{از حالت $s_1$:}
          \begin{itemize}
              \item با انتخاب اکشن $a_0$ به حالت پایانی ($s_3$) می‌رویم (پاداش = $10$).
              \item با انتخاب اکشن $a_1$ به حالت پایانی ($s_3$) می‌رویم (پاداش = $0$).
          \end{itemize}

    \item \textbf{از حالت $s_2$:}
          \begin{itemize}
              \item با انتخاب اکشن $a_0$ به حالت پایانی ($s_3$) می‌رویم (پاداش = $0$).
              \item با انتخاب اکشن $a_1$ به حالت پایانی ($s_3$) می‌رویم (پاداش = $10$).
          \end{itemize}
\end{itemize}

حال دو سیاست $\pi_1$ و $\pi_2$ را به این شکل تعریف می‌کنیم:
\begin{itemize}
    \item $\pi_1$: همواره در تمام حالت‌ها، اکشن $a_0$ را انتخاب می‌کند. در این حالت چون تمام اکشن‌ها و نتایج آنها قطعی هستند، می‌توانیم ارزش حالت‌ها را به سادگی مشخص کنیم:
          \begin{align*}
              V^{\pi_1}(s_0) = 10, \quad V^{\pi_1}(s_1) = 10, \quad V^{\pi_1}(s_2) = 0, \quad V^{\pi_1}(s_3) = 0
          \end{align*}
    \item $\pi_2$: همواره اکشن $a_1$ را انتخاب می‌کند. برای این سیاست داریم:
          \begin{align*}
              V^{\pi_2}(s_0) = 10, \quad V^{\pi_2}(s_1) = 0, \quad V^{\pi_2}(s_2) = 10, \quad V^{\pi_2}(s_3) = 0
          \end{align*}
\end{itemize}

سیاست $\pi_3 = \frac{1}{2}\pi_1 + \frac{1}{2}\pi_2$ را تعریف می‌کنیم (یعنی در هر گام با احتمال $\frac{1}{2}$ طبق سیاست $\pi_1$ و با احتمال $\frac{1}{2}$ مطابق $\pi_2$ عمل می‌کنیم). ارزش حالت‌ها در این سیاست را طبق این رابطه محاسبه می‌کنیم:
\begin{align*}
    V^{\pi_3}(s) = \sum_{a} \pi_3(a \mid s) \sum_{s'} T(s, a, s') \left[ R(s, a) + V^{\pi_3}(s') \right]
\end{align*}

بنابراین با توجه به این که $V^{\pi_3}(s_3) = 0$ است، داریم:
\begin{align*}
                & V^{\pi_3}(s_2) = \frac{1}{2}(0) + \frac{1}{2}(10) = 5 \\
                & V^{\pi_3}(s_1) = \frac{1}{2}(10) + \frac{1}{2}(0) = 5 \\
    \Rightarrow & V^{\pi_3}(s_0) = \frac{1}{2}(5) + \frac{1}{2}(5) = 5
\end{align*}

مشاهده می‌شود که $V^{\pi_3}(s_0) = 5$، اما طبق ادعای سوال باید مقدار آن $V^{\pi_3}(s_0) = \frac{1}{2}(10) + \frac{1}{2}(10) = 10$ باشد. بنابراین این ادعا نادرست است و برای آن مثال نقض آوردیم.